% Generated draft: requires mathematical and source review.
\documentclass[11pt]{article}
\usepackage[margin=25mm]{geometry}
\usepackage{fontspec}
\setmainfont{Latin Modern Roman}
\usepackage{amsmath,amssymb,mathtools,unicode-math}
\setmathfont{Latin Modern Math}
\usepackage{enumitem,needspace,xurl,hyperref}
\setlength{\parindent}{0pt}
\setlength{\parskip}{5pt}
\setlength{\emergencystretch}{3em}
\newcommand{\N}{\mathbb{N}}
\newcommand{\Z}{\mathbb{Z}}
\newcommand{\Q}{\mathbb{Q}}
\newcommand{\R}{\mathbb{R}}
\newcommand{\C}{\mathbb{C}}
\newcommand{\F}{\mathbb{F}}
\newcommand{\A}{\mathbb{A}}
\newcommand{\PP}{\mathbb{P}}
\newcommand{\E}{\mathbb{E}}
\newcommand{\eps}{\varepsilon}
\newcommand{\dd}{\,\mathrm{d}}
\begin{document}
\section*{UnsolvedMath ID: 20000464}
\textbf{Source problem code:} catalogue entry
\subsection{Definitions and mathematical statement}
The source asks whether harmonic analysis and nonabelian torsors give methods for finding rational points, and specifically whether a harmonic-analysis argument counts points on a quintic del Pezzo surface. A quintic del Pezzo surface is a smooth projective surface with ample \(-K_X\) and \(K_X^2=5\). The field, split/nonsplit form, counting open set and height are not specified. Wooley's comment that no combined examples were known is source-time context, not a newly verified literature assessment.

The attached aid uses an integral Cox-coordinate slice of the split quintic model. Its Cox equations have the Pl\"ucker form
\[
p_{ij}p_{kl}-p_{ik}p_{jl}+p_{il}p_{jk}=0
\quad(1\leq i<j<k<l\leq5),
\]
five equations in ten variables, together with the additional open/coprimality conditions needed for a torsor. It fixes four pairwise-coprime nonzero outer variables \(a_1,a_2,a_3,a_4\) and reports a bijective parametrization of the integral equation solutions by a rank-three congruence lattice
\[
L_a\subset\Z^3,\qquad\operatorname{covol}(L_a)=|a_1a_4|.
\]
The assignment of the outer variables, congruences and parametrization are not supplied, so this lattice is reported rather than explicitly reconstructed here.

Write \(L_a^*=\{\xi\in\R^3:\langle\xi,L_a\rangle\subset\Z\}\). For a smooth compactly supported weight \(w\), with \(\widehat w(\xi)=\int w(x)e^{-2\pi i\langle x,\xi\rangle}\,\dd x\), its reported exact fiberwise Poisson formula is
\[
\sum_{x\in L_a}w(x)=\frac1{|a_1a_4|}
\sum_{\xi\in L_a^*}\widehat w(\xi).
\]
The aid says its explicit dual parametrization uses two residue-frequency parameters, whose moduli and formulas are absent here. An exact formula on a fixed fiber does not supply a uniform estimate when the outer variables and divisibility conditions vary.

After the standard M\"obius inversions and height smoothing, the aid reduces a harmonic proof of the Manin--Peyre asymptotic to a uniform aggregate bound for nonzero dual frequencies and removal of smoothing. In the standard split rational model, for an anticanonical height and the open set obtained by removing accumulating exceptional curves, the intended target is
\[
N_{U,H}(B)\sim c_{X,H}B(\log B)^4.
\]
The missing bound must make the total nonzero-frequency contribution negligible at that scale, together with the necessary main-term and unsmoothing controls. Neither analytic bound is proved in the supplied aid. Its universal Cox torsor is a torus torsor; additive Poisson analysis of this slice is not a completed theory for arbitrary nonabelian torsors or general rational-point existence.
\subsection{Short English statement}
Seek harmonic-analysis methods combined with torsors, particularly for counting points on the degree-five del Pezzo surface. The attached aid reports a rank-three lattice parametrization and fiberwise Poisson formula, but leaves the global nonzero-frequency estimates and smoothing removal unproved.
\subsection{Sources}
\begin{itemize}
\item Local UnsolvedMath registry, permanent ID \(20000464\), exact public code \texttt{catalogue entry}; canonical location \texttt{aim-algebraic-number-theory-notes.json}, notes[115].
\item AIM workshop ``Rational and integral points on higher dimensional varieties,'' Question 31; supplied URL \url{https://aimath.org/WWN/qptsurface2/qptsurface2.pdf}. Dialogue is attributed to Takloo-Bighash, Hassett and Wooley. The verification removes the trailing ``Problem/'' as an extraction/navigation artifact. PDF/HTML were not consulted here; the literature field is empty.
\item Attached research attempt 1, ``A fiberwise harmonic bridge on the quintic del Pezzo torsor,'' referring to Cox/Pl\"ucker models, Poisson summation and the Manin--Peyre asymptotic without complete bibliographic citations.
\end{itemize}
\end{document}
